% fig:adni-pca: both directions' flow states in a PCA fitted on the endpoint codes only.
% Linear projection on purpose -- relative distance and straightness survive it, which is
% not true of t-SNE or UMAP. Black = mean trajectory numbered by block, grey = ~25
% individual slices, clouds = the endpoint codes. Drawn by inpaper_utils/domain_probe.py
% from the pooled features of inpaper_utils/latent_geometry.py; the quantitative claims in
% the text (path/chord, probe margin) are measured in the latent space, not in this
% projection.
\begin{figure}[t]
\centering
\includegraphics[width=\linewidth]{figures/adni_pca_path.pdf}
\caption{The states of both directions over the $430$ held-out slices, projected onto a
PCA fitted on the endpoint codes $E_{\mathrm{T1}}(x)$ and $E_{\mathrm{FA}}(y)$ alone:
(a) T1$\rightarrow$FA and (b) FA$\rightarrow$T1. Black is the mean trajectory, numbered by
block; thin grey lines are individual slices; the two clouds are the endpoint codes
themselves. The projection is linear, so relative distance and straightness survive it,
which is not true of a $t$-SNE or UMAP embedding; the numbers quoted in the text are
measured in the latent space itself, not in this projection.}
\label{fig:adni-pca}
\end{figure}
